% Chapter introduction: why the topic arose, who first asked, what it answers.
% STATUS: written 2026-11-05 (all four body sections sourced and gate-clean).
% The introduction reports the four threads of the chapter and cites only the
% sources that carry them (Freudenburg 2017, Drensky--Makar-Limanov 2009,
% Baltazar 2016, Rewri--Kour 2025, Dixmier 1968, Belov-Kanel--Kontsevich 2007,
% Hazewinkel 2013, Sweedler 1980, Nowicki 2004, Lequain 2008, Ore 1933).
% The classical first sources are named through the papers that quote them.
\emph{The derivative returns to its first home.} Chapter~\ref{ch:origins}
opened the book with Leibniz and the derivative of a polynomial. This
chapter comes back to that home and asks what the Leibniz rule becomes on
the polynomial ring $k[x_{1},\dots,x_{n}]$ and on the noncommutative rings
that imitate it. Four threads answer, each read from the papers that carry
it.

\hbfigure{12-polynomial-weyl/weyl-ladder}{The Weyl algebra acting on $k[x]$: multiplication by $x$ raises the degree; $\partial = d/dx$ lowers it, with $\partial(x^{n}) = n x^{n-1}$; and $\partial x - x\partial = 1$.}{fig:weyl:ladder}

The first thread asks what a derivation of a polynomial ring can be
\cite[p.~75]{Freudenburg2017}. the fourteenth problem of Hilbert and the
Jacobian conjecture made the question urgent
\cite[p.~xii]{Freudenburg2017}. Weitzenb\"ock gave the first answer in
1932: for a linear locally nilpotent derivation the kernel is finitely
generated \cite[p.~41]{DrenskyMakarLimanov2009}.
Section~\ref{sec:polynomial-weyl:polynomial} defines the named families:
locally nilpotent, triangular, elementary, monomial, Weitzenb\"ock
\cite[p.~76]{Freudenburg2017}, Shamsuddin \cite[p.~2033]{Baltazar2016} and
Lotka--Volterra \cite[p.~2]{RewriKour2025}. It then reads the kernel
theorems. These are the Weitzenb\"ock theorem and the confirmed Nowicki
conjecture \cite[Theorem~3, p.~44]{DrenskyMakarLimanov2009}. The section
closes with the fixed-group theorems of Mendes and Pan
\cite[Theorem~1, p.~2]{MendesPan2017} and of Rewri and Kour
\cite[Theorem~4.1, p.~21]{RewriKour2025}. The automorphisms that commute
with a derivation $d$ form the fixed group of $d$ there. Bedratyuk gives an
algorithm for the kernel of a Weitzenb\"ock derivation, built from Casimir
elements of the derivation \cite[p.~1]{Bedratyuk2005Casimir}.
Mishra, Mondal and Shukla compute the fixed group of a monomial derivation
whose exponent matrix is invertible \cite[Theorem~3.1,
p.~5]{MishraMondalShukla2026Monomial}. On a simple derivation that takes no
unit value, the fixed group reduces to the translations
\cite[Theorem~2.4, p.~5]{MishraMondalShukla2025Simple}. A derivation of a
polynomial ring is a first-order step of a Hasse--Schmidt system, and the
four named families live on the polynomial ring as such steps. Locally
nilpotent derivations are the nil case: each polynomial dies under some
power of the derivation \cite[p.~3]{Freudenburg2017}. The general
definitions live in
Section~\ref{sec:higher-multi:hasse-schmidt} and
Section~\ref{sec:basic-theory:definitions}.

The second thread carries the derivative into a noncommutative algebra
\cite[p.~209]{Dixmier1968}. Dixmier named the Weyl algebra in 1968, after
a suggestion of Segal, and listed six problems at the end of his paper
\cite[p.~241]{Dixmier1968}. Section~\ref{sec:polynomial-weyl:weyl} reads
his definition, his structure theorem for $A_{n}$ and his first problem,
now the Dixmier conjecture. Belov-Kanel and Kontsevich prove that this
conjecture agrees, in a stable sense, with the Jacobian conjecture
\cite[Theorem~1, p.~2]{BelovKanelKontsevich2007}; the two threads of the
chapter close on the same question.

The third thread counts the order of an operator
\cite[p.~4]{Hazewinkel2013}. Grothendieck defined the differential
operators of a commutative algebra by commuting them with multiplications;
Section~\ref{sec:polynomial-weyl:differential-operators} recalls the
recursive characterization as Hazewinkel states it
\cite[Theorem~6.2, p.~4]{Hazewinkel2013}. Sweedler carried the same
commutator to noncommutative algebras and named the right derivation
\cite[p.~29]{Sweedler1980}, the first kind of the chapter with the right
hand rule of Section~\ref{sec:basic-theory:definitions}.

The fourth thread builds polynomials over a derivation
\cite[p.~480]{Ore1933}.
Section~\ref{sec:polynomial-weyl:ore-differential}
takes the differential form $R[t;\delta]$ of the Ore extension of
Section~\ref{sec:generalized-twisted:ore-extensions} to its concrete
ground.
The Weyl algebra is the first example, and the $d$-simple rings make
$R[t;\delta]$ simple \cite[p.~801]{Lequain2008}.
On $k[x][t;\delta]$ every local derivation is already a derivation
\cite[Theorem~4.1, p.~4569]{Nowicki2004Local}.
Chapter~\ref{ch:commutative}
follows the derivations of polynomial rings into geometry.
